\SourceLeaf{522}Hence $\Sigma[-\mathfrak{Q}]\times\Sigma[\mathfrak{Q}]$ expands into
\[
\begin{aligned}
&-[0]+[1]+[4]+[9]\dotsfive+[(n-1)^2]\\
&+[0]+[3]+[8]+[15]\dotsfive+[nn-1]\\
&+[0]+[5]+[12]+[21]\dotsfive+[nn+2n-3]\\
&+[0]+[7]+[16]+[27]\dotsfive+[nn+4n-5]\\
&+\text{etc.}\\
&+[0]+[2n-1]+[4n]+[6n+3]\dotsfive+[3nn-6n+3]
\end{aligned}
\]
Summing these parts \textit{vertically} gives
\[
\begin{aligned}
&n[0]+[1]\times\frac{1-[2n]}{1-[2]}+[4]\frac{1-[4n]}{1-[4]}+[9]\frac{1-[6n]}{1-[6]}\\
&\qquad+\text{etc.}+[(n-1)^2]\times\frac{1-[2nn-2n]}{1-[2n-2]}
\end{aligned}
\]
In this expression all parts except the first vanish whenever $n$ is odd; for then all of $1-[2n]$, $1-[4n]$, $1-[6n]$, etc. become $=0$, whereas none of the denominators $1-[2]$, $1-[4]$, $1-[6]$, $1-[8]$, etc. up to $1-[2n-2]$ does. When $n$ is even, however, one of the denominators is also $=0$, namely $1-[n]$, to which the term $[\frac14nn]\times\frac{1-[nn]}{1-[n]}$ corresponds; the sum of the parts from which this term arose becomes $=n[\frac14nn]$. Here again we must distinguish two cases. When $n$ is divisible by four, $\frac14nn\equiv0\, (\mathrm{mod}.n)$ and therefore $[\frac14nn]=1$; but when $n$ is twice an odd number, $\frac14nn\equiv\frac12n\, (\mathrm{mod}.n)$ and therefore necessarily $[\frac14nn]=-1$. We finally conclude that
\par\noindent 1) for an odd value of $n$, the required product is $=n$
\par\noindent 2) for a value divisible by four, the product is $=2n$
\par\noindent 3) for a value that is twice an odd number, it is $=0$.\quad Q. E. I.
\par\medskip\Needspace{4\baselineskip}\pdfbookmark[1]{371.}{sectio-octava-art-371}\centerline{\textbf{371.}}\nobreak\smallskip
It is now worth considering the nature of the sum $\Sigma[\mathfrak{Q}]$ more closely.
\par I. Since we may replace the squares $0$, $1$, $4$, $9$, $16$, etc. by their least residues modulo $n$, it is evident that if $M$ denotes, without further specification, the quadratic residues of the number $n$ from $0$ to $n-1$, and $m$ the number of roots of the congruence $xx\equiv M(\mathrm{mod}.m)$, then $\Sigma[\mathfrak{Q}]=\Sigma m[M]$. We showed how to determine the number $m$ in Articles $104$, $105$.
\par II. If $n$ is an (odd) prime, then for $M=0$ we have $m=1$, whereas for every other value of $M$ we have $m=2$. If $n$ is a power of an odd prime $=p^\nu$, then $m=2$ for every value of $M$ not divisible by $p$\quad\textemdash
\par\nobreak\noindent\hbox to\linewidth{\textemdash\hfil\textemdash\hfil\textemdash\hfil\textemdash\hfil\textemdash\hfil\textemdash\hfil\textemdash\hfil\textemdash\hfil\textemdash\hfil\textemdash\hfil\textemdash\hfil\textemdash\hfil\textemdash\hfil\textemdash\hfil\textemdash\hfil\textemdash\hfil\textemdash}\par
